\chapter{Related Work}\label{chap:related_work}

\section{Subjective \acs{IQA}}\label{sec:subjective_iqa}

Subjective image quality assessment is grounded in human visual perception and remains the gold standard for evaluating visual fidelity. Standardized by \acs{ITU-R BT.500-15}~\cite{itu2023r}, these methodologies rely on controlled laboratory conditions, calibrated displays, and predefined rating procedures to ensure reproducibility and validity.

\subsection{Assessment Protocols}\label{sec:assessment_protocols}

Protocols are typically divided into two categories: quality assessment, which captures the overall perceived quality of an image, and impairment assessment, which measures degradation relative to a reference. Two methodological paradigms are most common: the \acf{SS} method, where each image is rated independently, and the \acf{DS} method, where reference and test images are evaluated comparatively. SS protocols dominate modern datasets due to their simplicity, efficiency, and reduced observer fatigue.

Rating scales can be numerical or categorical. Numerical scales range from 1 to 5, 1 to 11, or 0 to 100, and are used to express either quality (excellent to bad) or impairment (imperceptible to very annoying). Categorical scales rely on descriptive labels, and are commonly found in ITU protocols. The choice of scale depends on whether the target measure is \acf{MOS} or \acf{DMOS}, and on the desired sensitivity of the test.

Each test session consists of multiple presentations. For a given presentation defined by condition $j$, image or sequence $k$, and repetition $r$, observer $i$ provides a discrete score $u_{ijkr}$. The MOS is computed as:

\begin{equation}
\text{MOS}_{jkr} = \frac{1}{N} \sum_{i=1}^{N} u_{ijkr}
\end{equation}
where $N$ is the number of observers. In double-stimulus protocols, the DMOS is defined as the difference between the MOS of the reference and that of the distorted image:

\begin{equation}
\text{DMOS}_{jkr} = \text{MOS}_{\text{ref},jkr} - \text{MOS}_{\text{dist},jkr}
\end{equation}

These values are often accompanied by confidence intervals to quantify statistical uncertainty and by post-screening to remove inconsistent observers. The specific derivations, including standard deviation, confidence intervals, and the ITU-recommended kurtosis-based rejection procedure, are detailed in Appendix~\ref{app:subjective-screening}.

\subsection{Crowdsourcing and Remote Studies}\label{sec:crowdsourcing}

Laboratory tests provide high internal validity but are costly, time-consuming, and limited in participant diversity. To overcome these constraints, recent works have turned to crowdsourcing platforms such as Amazon Mechanical Turk, Figure Eight, and Prolific~\cite{ghadiyaram2015massive, koniq10k}. At web scale, researchers obtain large, diverse panels and ratings gathered under everyday viewing settings, which improves ecological validity for images consumed on heterogeneous devices.

However, uncontrolled environments introduce noise: participants use uncalibrated displays, ratings may be inconsistent, and some responses are malicious or inattentive. To mitigate these issues, robust post-screening strategies are employed, including consistency checks, gold-standard reference images, and statistical filtering~\cite{ghadiyaram2015massive, koniq10k}. Despite these limitations, crowdsourced \acs{IQA} datasets have proven critical for training modern deep learning-based models, as they provide both scale and diversity unavailable in controlled laboratory conditions. These trade-offs shape how \acs{IQA} datasets are constructed and interpreted.

\subsection{Benchmark Datasets}\label{sec:benchmark_datasets}

Several benchmark datasets have been curated to support the development and evaluation of \acs{IQA} models. Early datasets such as LIVE~\cite{livewebsite, sheikh2005live, sheikh2006live}, CSIQ~\cite{csiq}, TID2008~\cite{tid2008}, and TID2013~\cite{tid2013} follow \acs{ITU-R BT.500-15}-compliant protocols, providing high-quality \acs{MOS} or \acs{DMOS} annotations across a wide variety of synthetic distortion types. These databases remain the most widely used for evaluating traditional full-reference \acs{IQA} metrics.

The IVC~\cite{ivc} dataset further contributed subjective quality scores for compression-related distortions, while SIQAD~\cite{siqad} targeted screen content images, extending \acs{IQA} research to non-natural imagery. KADID-10k~\cite{lin2019kadid} scaled this approach to 10,000 distorted images with 25 distortion types and multiple intensity levels, enabling large-scale training and fine-grained benchmarking. The Waterloo IAA database~\cite{waterloo} emphasized robustness testing through massive-scale synthetic distortions and helped reveal model overfitting issues to specific artifact distributions.

More recent contributions shifted toward naturalistic and crowdsourced assessments. KonIQ-10k~\cite{koniq10k} introduced 10,073 images collected from the web with authentic distortions and over 1.2 million quality ratings, making it one of the largest ecologically valid \acs{IQA} resources. SPAQ~\cite{spaq} specifically addressed smartphone photography, collecting 11,125 images and crowdsourced perceptual scores to capture the variability of mobile imaging pipelines. PIPAL~\cite{pipal} targeted perceptual quality in image restoration, including outputs from generative models, and has become a key benchmark for deep learning-based perceptual metrics. Similarly, BAPPS~\cite{bapps} provided a large-scale pairwise comparison dataset for learning perceptual similarity directly from human judgments.

Beyond general-purpose \acs{IQA}, domain-specific datasets remain critical. Remote sensing benchmarks such as ROSIS~\cite{rosis}, AVIRIS~\cite{aviris}, and HYDICE~\cite{hydice} provide hyperspectral imagery where distortions affect both perceptual and functional quality, but they lack standardized \acs{MOS} or \acs{DMOS} annotations. These datasets illustrate the difficulty of extending \acs{IQA} methodologies into application-driven contexts.

The \acf{FRLL}~\cite{frll} represents a distinct contribution to IQA benchmarking, focusing on high-quality facial \acs{ICAO} compliant imagery. It contains 102 unique subjects photographed under controlled studio conditions, with consistent illumination, wearing plain white t-shirts, and neutral, frontal poses. In addition to meeting passport-style requirements, the dataset also includes self-reported age, gender, and ethnicity. Attractiveness ratings (on a 1-7 scale from “much less attractive than average” to “much more attractive than average”) for the neutral front faces, provided by 2513 raters aged 17-90, are also included.

This combination of demographic attributes and attractiveness ratings enables not only research that bridges perceptual quality with subjective aesthetic judgments, but also systematic analysis of demographic bias in human evaluations of facial images. These dual annotations make the dataset particularly suitable for studies involving biometric verification and \acs{MRTDs}, where both technical quality and fairness considerations are critical.

Unlike general-purpose IQA benchmarks, \acs{FRLL} does not include traditional synthetic distortions. Instead, its value lies in providing a pristine reference set of demographically diverse facial images, which can be systematically altered or embedded with task-driven degradations such as steganography, compression, or adversarial perturbations. This makes \acs{FRLL} an essential testbed for application-driven IQA, where the objective extends beyond generic perceptual fidelity to domain-specific measures of utility and fairness.

Together, these databases trace the evolution of \acs{IQA} benchmarking from controlled laboratory studies with synthetic distortions to large-scale, crowdsourced, and application-specific datasets. This progression not only enabled the development of modern no-reference and deep learning-based \acs{IQA} models but also revealed persistent challenges in generalization to real-world, domain-critical scenarios.

\section{Objective \acs{IQA}}\label{sec:objective_iqa}

Objective \acs{IQA} refers to the automatic estimation of visual quality using computational models. Unlike subjective approaches that rely on human ratings, objective \acs{IQA} enables consistent and scalable evaluations, making it essential for applications such as compression, transmission, restoration, and enhancement~\cite{gonzalez2009digital, moorthy2011blind, lin2020surpassing}. Over the past decades, numerous methods have been proposed, shaped by different assumptions about the \acf{HVS} and the statistical properties of natural images. Surveys and comparative studies~\cite{shahrukh2019survey, wang2004image, liu2012fusion, ding2020dists} summarize this evolution, showing a progression from simple error-based fidelity measures to advanced perceptual and learning-based frameworks. Although no metric fully replicates human judgment, objective \acs{IQA} provides a reproducible and low-cost alternative that supports large-scale benchmarking and real-time quality monitoring.

A principled way to classify \acs{IQA} methods is through two dimensions: the computational domain in which comparisons are performed and the spectral representation adopted.

\subsection{Computational Domains}

The computational domain defines the signal space where similarity is quantified.

\begin{itemize}
  \item Spatial domain: Early metrics such as MSE, RMSE, and PSNR operate directly on pixel intensities by measuring absolute differences or signal-to-noise ratios~\cite{gonzalez2009digital}. Despite their simplicity, they fail to account for perceptual masking and structural consistency.
  \item Frequency domain: Transform-based methods use the discrete cosine transform (DCT), discrete Fourier transform (DFT), or wavelets to capture energy distributions across spatial frequencies. Examples include IFC~\cite{sheikh2005ifc} and VSNR~\cite{chandler2007vsnr}, which exploit properties of the \acs{HVS} such as contrast sensitivity.
  \item Gradient and saliency domains: These methods emphasize structural edges and perceptually important regions. The gradient magnitude similarity deviation (GMSD)~\cite{xue2014gmsd} measures gradient consistency, while saliency-based models~\cite{hou2007saliency} weight distortions according to visual attention maps.
  \item Information-theoretic domain: Some metrics interpret quality as the amount of shared information between reference and distorted images. VIF~\cite{sheikh2006vif} and related approaches build on mutual information, quantifying how much perceptual information can be extracted reliably from a degraded signal.
  \item Deep feature domain: More recent approaches leverage convolutional neural networks or vision transformers to learn perceptual embeddings. Metrics such as LPIPS~\cite{zhang2018lpips} and DISTS~\cite{ding2020dists} correlate strongly with human judgments by comparing features extracted from pre-trained networks.
\end{itemize}

\subsection{Spectral Representations}

The spectral representation specifies the channel configuration used for quality computation.

\begin{itemize}
  \item Luminance-based: Many classical approaches, such as SSIM~\cite{wang2004image}, focus solely on the luminance channel (Y). This reflects the dominant role of brightness perception in the \acs{HVS} and simplifies computation.
  \item Color-based: Extensions to RGB or alternative color spaces, including YCbCr and CIELab, enable the incorporation of chromatic information. Examples include FSIM and its color extension FSIMc~\cite{zhang2011fsim}, or PQM~\cite{wang2002uqi}, which attempt to capture hue and saturation degradations alongside luminance.
  \item Multispectral and hyperspectral: In domains such as remote sensing, quality measures generalize to dozens or even hundreds of bands. Metrics like SAM~\cite{yuhas1992sam}, ERGAS~\cite{wald2000ergas}, and $D_\lambda$~\cite{ivc} ensure spectral consistency across wavelengths and are critical for preserving material properties.
\end{itemize}

This dual perspective highlights the diversity of \acs{IQA} methodologies and provides a framework for their categorization. Building on it, Table~\ref{tab:fr_iqa_metrics} presents the 40 \acs{FR} metrics and Table~\ref{tab:nr_iqa_metrics} summarizes the 4 \acs{NR} metrics considered in this work. Each metric is characterized by its reference type, year of introduction, category domain, and tested databases, showing how the field has advanced from simple pixel-wise error measures to perceptually and semantically informed quality predictors.

\subsection{Full-Reference \acs{IQA}}\label{sec:full_reference_iqa}

\begin{longtable}{
p{0.13\linewidth}  % Metric
p{0.05\linewidth}  % Year
p{0.22\linewidth}  % Category
p{0.50\linewidth}  % Tested Databases
}
    \caption{Summary of the 40 \acs{FR}-\acs{IQA} metrics used}\label{tab:fr_iqa_metrics}\\
    \hline % chktex 44
    \textbf{Metric} & \textbf{Year} & \textbf{Category} & \textbf{Tested Databases} \\
    \hline % chktex 44
    \endfirsthead % chktex 1
    \multicolumn{4}{r}{\small\itshape(continuation)}\\ % chktex 6
    \hline % chktex 44
    \textbf{Metric} & \textbf{Year} & \textbf{Category} & \textbf{Tested Databases} \\
    \hline % chktex 44
    \endhead % chktex 1
    \hline % chktex 44
    \endfoot % chktex 1

    MSE~\cite{wang2004image} & 2004 & Error-based & LIVE, CSIQ, TID2008, TID2013, KADID-10k \\ \hline
    RMSE~\cite{wang2004image} & 2004 & Error-based & LIVE, CSIQ, TID2008, TID2013, KADID-10k \\ \hline
    MAE~\cite{wang2004image} & 2004 & Error-based & LIVE, CSIQ, TID2008, TID2013, KADID-10k \\ \hline
    PSNR~\cite{wang2004image} & 2004 & Error-based & LIVE, CSIQ, TID2008, TID2013, KADID-10k \\ \hline
    PSNR-B~\cite{ma2011psnrb} & 2011 & Error-based & LIVE, CSIQ, TID2013 \\ \hline
    SNR~\cite{wang2004image} & 2004 & Error-based & LIVE, CSIQ, TID2008 \\ \hline
    SAM~\cite{yuhas1992sam} & 1992 & Spectral-based & AVIRIS, HYDICE, ROSIS, TID2013 \\ \hline
    ERGAS~\cite{wald2000ergas} & 2000 & Spectral-based & AVIRIS, ROSIS, TID2013, KADID-10k \\ \hline
    $D_\lambda$~\cite{wald2000ergas} & 2000 & Spectral-based & AVIRIS, ROSIS, TID2013, KADID-10k \\ \hline
    SSIM~\cite{wang2004image} & 2004 & Structural Similarity & LIVE, CSIQ, TID2008, TID2013 \\ \hline
    DSSIM~\cite{wang2004image} & 2004 & Structural Similarity & LIVE, CSIQ, TID2008 \\ \hline
    MS-SSIM~\cite{wang2003multiscale} & 2003 & Structural Similarity & LIVE, CSIQ, TID2013, KADID-10k \\ \hline
    IW-SSIM~\cite{li2011iwssim} & 2010 & Structural Similarity & LIVE, CSIQ, TID2013 \\ \hline
    L-SSIM~\cite{wang2004image} & 2004 & Structural Similarity & --- \\ \hline
    UQI~\cite{wang2002uqi} & 2002 & Structural Similarity & --- \\ \hline
    SCC~\cite{eskicioglu1995scc} & 1995 & Structural Similarity & --- \\ \hline
    CW-SSIM~\cite{sampat2009cwssim} & 2009 & Wavelet Similarity & LIVE, CSIQ, TID2013 \\ \hline
    W-SSIM~\cite{li2009wssim} & 2009 & Wavelet Similarity & LIVE \\ \hline
    ESSIM~\cite{chen2006esim} & 2006 & Edge Similarity & LIVE \\ \hline
    G-SSIM~\cite{chen2006gssim} & 2006 & Gradient Similarity & LIVE, TID2008 \\ \hline
    R-SSIM~\cite{yan2015rssim} & 2015 & Gradient Similarity & --- \\ \hline
    GMS~\cite{xue2014gmsd} & 2014 & Gradient Similarity & LIVE, CSIQ, TID2008 \\ \hline
    GMSD~\cite{xue2014gmsd} & 2014 & Gradient Similarity & LIVE, CSIQ, TID2008 \\ \hline
    C-SSIM~\cite{hassan2017cssim} & 2017 & Color Similarity & TID2013 \\ \hline
    FSIM~\cite{zhang2011fsim} & 2011 & Feature Similarity & LIVE, CSIQ, TID2008, TID2013, KADID-10k \\ \hline
    FSIMc~\cite{zhang2011fsim} & 2011 & Feature Similarity & LIVE, CSIQ, TID2008, IVC \\ \hline
    MS-FSIM~\cite{zhang2011fsim} & 2011 & Color Similarity & --- \\ \hline
    SR-SIM~\cite{zhang2012srsim} & 2013 & Feature Similarity & LIVE, CSIQ, TID2008 \\ \hline
    VSNR~\cite{chandler2007vsnr} & 2007 & Feature Similarity & LIVE, CSIQ \\ \hline
    IFC~\cite{sheikh2005ifc} & 2005 & Information & LIVE \\ \hline
    VIF~\cite{sheikh2006vif} & 2006 & Information & LIVE, CSIQ, TID2008 \\ \hline
    VIFp~\cite{sheikh2006vif} & 2006 & Information & LIVE \\ \hline
    VIFc~\cite{sheikh2006vif} & 2006 & Information & --- \\ \hline
    LPIPS-Alex~\cite{zhang2018lpips} & 2018 & Deep Learning & BAPPS, LIVE, CSIQ, TID2013 \\ \hline
    LPIPS-Squeeze~\cite{zhang2018lpips} & 2018 & Deep Learning & BAPPS, LIVE, CSIQ, TID2013 \\ \hline
    LPIPS-VGG~\cite{zhang2018lpips} & 2018 & Deep Learning & BAPPS, LIVE, CSIQ, TID2013 \\ \hline
    DISTS~\cite{ding2020dists} & 2020 & Deep Learning & LIVE, CSIQ, TID2013, KADID-10k \\ \hline
    WaDIQaM~\cite{bosse2017wadiqam} & 2017 & Deep Learning & LIVE, TID2013 \\ \hline
    PSM & --- & --- & --- \\ \hline
    SPIQ & --- & --- & --- \\ \hline
\end{longtable}

% =========================
% RELATED WORK: from metrics tables downward
% =========================

\section{Objective IQA: critical comparison and scope}\label{sec:objective_iqa_critical}
This section contrasts major metric families used in this thesis and motivates which ones are retained for experiments. Detailed formulas and notations are deferred to Appendix~\ref{appendix:metrics} to keep the narrative focused.

\subsection{Families of full-reference metrics}\label{sec:fr_families}
We group FR metrics into four families: error-based, similarity-based, information-theoretic, and deep learning feature-based. For each family we summarize principle, strengths, failure modes, and expected behavior on subtle distortions.

\subsubsection{Error-based}
Principle: pixel-wise deviations.
Pros: simple, fast, interpretable.
Cons: weak perceptual alignment; insensitive to structural rearrangements; often saturate on subtle distortions.
Use in this thesis: baseline correlation analysis and fusion inputs. See Appendix~\ref{appendix:metrics_error} for formulas.

\subsubsection{Similarity-based}
Principle: structural or feature resemblance in luminance, contrast, texture, phase.
Pros: improved alignment with human judgments; robust to some luminance shifts.
Cons: sensitivity to geometric misalignment; some variants overfit to specific artifacts.
Use in this thesis: primary contributors to fusion. See Appendix~\ref{appendix:metrics_similarity}.

\subsubsection{Information-theoretic}
Principle: mutual information preservation under natural scene statistics models.
Pros: principled; captures information loss.
Cons: higher complexity; modeling assumptions can misfit faces or non-natural images.
Use in this thesis: complementary signals in fusion. See Appendix~\ref{appendix:metrics_information}.

\subsubsection{Deep learning feature-based}
Principle: distances in pretrained deep feature spaces; structure and texture factors.
Pros: high correlation with MOS; captures complex distortions.
Cons: data bias risk; less interpretable; backbone choice matters.
Use in this thesis: tested as fusion inputs and compared in correlation. See Appendix~\ref{appendix:metrics_deep}.

\subsection{FR metrics considered in this thesis}\label{sec:fr_metrics_table}
We report the FR metrics actually computed. Equations are in Appendix~\ref{appendix:metrics}.

\begin{longtable}{
p{0.18\linewidth}  % Metric
p{0.07\linewidth}  % Year
p{0.23\linewidth}  % Family/Category
p{0.42\linewidth}  % Representative tested databases
}
\caption{Full-reference IQA metrics considered in this thesis. Formulas in Appendix~\ref{appendix:metrics}.}\label{tab:fr_iqa_metrics_compact}\\
\hline
Metric & Year & Family/Category & Representative databases \\
\hline
\endfirsthead
\multicolumn{4}{r}{(continuation)}\\
\hline
Metric & Year & Family/Category & Representative databases \\
\hline
\endhead
\hline
\endfoot

% TODO: keep only metrics you actually used; keep rows concise
MSE & 2004 & Error-based & LIVE, CSIQ, TID2013 \\
RMSE & 2004 & Error-based & LIVE, CSIQ, TID2013 \\
MAE & 2004 & Error-based & LIVE, CSIQ \\
PSNR & 2004 & Error-based & LIVE, CSIQ, TID2013, KADID-10k \\
PSNR-B & 2011 & Error-based (blocking aware) & LIVE, CSIQ, TID2013 \\
SNR & 2004 & Error-based & LIVE, CSIQ \\
SSIM & 2004 & Similarity-based & LIVE, CSIQ, TID2008, TID2013 \\
MS-SSIM & 2003 & Similarity-based (multiscale) & LIVE, CSIQ, TID2013, KADID-10k \\
IW-SSIM & 2010 & Similarity-based (info-weighted) & LIVE, CSIQ, TID2013 \\
CW-SSIM & 2009 & Similarity-based (wavelet phase) & LIVE, CSIQ \\
W-SSIM & 2009 & Similarity-based (region-weighted) & LIVE \\
ESSIM & 2006 & Similarity-based (edge) & LIVE \\
G-SSIM & 2006 & Similarity-based (gradient) & LIVE, TID2008 \\
R-SSIM & 2015 & Similarity-based (region) & reported in literature \\
L-SSIM & 2004 & Similarity-based (local) & reported in literature \\
UQI & 2002 & Similarity-based (early model) & reported in literature \\
SCC & 1995 & Similarity-based (correlation) & reported in literature \\
FSIM & 2011 & Feature similarity (PC+GM) & LIVE, CSIQ, TID2008, TID2013 \\
FSIMc & 2011 & Feature similarity with color & LIVE, CSIQ, TID2008, IVC \\
SR-SIM & 2013 & Saliency-weighted similarity & LIVE, CSIQ, TID2008 \\
VSNR & 2007 & HVS thresholding, frequency & LIVE, CSIQ \\
IFC & 2005 & Information-theoretic & LIVE \\
VIF & 2006 & Information-theoretic & LIVE, CSIQ, TID2008 \\
VIFp & 2006 & Information-theoretic (pixel) & LIVE \\
VIFc & 2006 & Information-theoretic (color) & reported in literature \\
LPIPS-Alex & 2018 & Deep features & BAPPS, LIVE, CSIQ, TID2013 \\
LPIPS-Squeeze & 2018 & Deep features & BAPPS, LIVE, CSIQ, TID2013 \\
LPIPS-VGG & 2018 & Deep features & BAPPS, LIVE, CSIQ, TID2013 \\
DISTS & 2020 & Deep features (structure+texture) & LIVE, CSIQ, TID2013, KADID-10k \\
WaDIQaM & 2017 & Deep learning FR (siamese) & LIVE, TID2013 \\
% Add any additional FR metrics you actually computed here
\hline
\end{longtable}

\subsection{No-reference metrics included for comparison}\label{sec:nr_metrics_table}
We only list NR metrics that appear later in experiments. Formulas in Appendix~\ref{appendix:metrics_nr}.

\begin{longtable}{
p{0.18\linewidth}
p{0.07\linewidth}
p{0.23\linewidth}
p{0.42\linewidth}
}
\caption{No-reference IQA or FIQA baselines included in this thesis.}\label{tab:nr_iqa_metrics_compact}\\
\hline
Metric & Year & Type & Representative datasets \\
\hline
\endfirsthead
\multicolumn{4}{r}{(continuation)}\\
\hline
Metric & Year & Type & Representative datasets \\
\hline
\endhead
\hline
\endfoot

PIQE & 2016 & Classical NR-IQA & LIVE, CSIQ, TID2013 \\
NIQE & 2013 & Classical NR-IQA (NSS) & LIVE, CSIQ, TID2013 \\
SER-FIQ & 2020 & NR-FIQA (stochastic embedding) & face benchmarks \\
MagFace & 2021 & FIQA (quality-aware embedding) & face benchmarks \\
\hline
\end{longtable}

\subsection{What each family gets right and wrong on subtle distortions}\label{sec:family_tradeoffs}
Error-based metrics underreact to imperceptible yet task-relevant changes. Similarity-based metrics improve alignment but can be brittle to alignment issues. Information-theoretic metrics capture information loss but depend on modeling assumptions. Deep feature metrics capture complex changes but inherit bias from pretraining. This motivates metric fusion to combine complementary signals.

\section{Fusion of FR metrics and learning from pseudo-labels}\label{sec:fusion_related}
Prior work has fused heterogeneous metrics and exploited pseudo-labels to overcome limited MOS. We place our pipeline in this line of research and point to Appendix~\ref{appendix:regressors} for model details.

\subsection{Regression models commonly used for fusion}\label{sec:regression_models_overview}
We briefly list models evaluated in this thesis for fusion:
Linear and regularized linear models, kernel SVR, tree ensembles, gradient boosting families, and categorical-boosting approaches.
Rationales: linear models for interpretability, SVR for flexible margins on small data, ensembles for non-linear interactions, boosting for strong tabular performance.

\subsection{Why gradient boosting and why CatBoost here}\label{sec:why_catboost}
CatBoost provides ordered boosting and robust regularization on small to medium tabular data, captures non-linearities among metrics, and supports feature importance analysis. Details of CatBoost configuration and training used in this thesis are presented in Methodology, Section~\ref{sec:fusion_fr_to_pmos} and Section~\ref{sec:catboost_details}. Comparative descriptions of alternative regressors are in Appendix~\ref{appendix:regressors}.

\section{Steganography and printer-proof approaches}\label{sec:steganography_related}
We review steganography with emphasis on methods robust to print-scan pipelines, which is the operational context for MRTDs.

\subsection{Classical embedding domains}\label{sec:steg_classical_domains}
Spatial LSB, transform-domain DCT and DWT, and adaptive schemes. Strengths, detectability, and robustness trade-offs.

\subsection{Printer-proof deep pipelines}\label{sec:steg_printer_proof}
End-to-end learned encoders and decoders trained with differentiable print-scan models and geometric-photometric perturbations.

\subsubsection{StegaStamp}
Key idea, robustness sources, message capacity, typical artifacts. Relevance to high-resolution face portraits.

\subsubsection{CodeFace}
Face-specific alignment and identity-preserving losses. Relevance to FIQA and MRTD portraits.

\subsubsection{RiemStega}
Statistical descriptors on SPD manifolds to preserve global structure. Expected behavior under scanner noise and color drift.

\subsubsection{StampOne}
Frequency-aware encoding with attention across bands. Robustness to printer color shifts and sensor noise.

\section{Synthesis and implications for our approach}\label{sec:related_synthesis}
Takeaways. No single FR metric captures imperceptible but recognition-relevant degradations. Families are complementary but biased in distinct ways. Pseudo-label fusion offers a scalable path to align with MOS. Printer-proof steganography demands quality predictors sensitive to subtle, structured perturbations. These observations motivate the fusion-to-NR strategy adopted in this thesis and lead directly to the pipeline in Chapter~\ref{chap:methodology}.

% =========================
% APPENDICES: referenced by the skeleton above
% =========================

\appendix

\chapter{Mathematical definitions of IQA metrics}\label{appendix:metrics}
This appendix collects formulas and notations for all metrics referenced in Chapter~\ref{chap:related_work}. Keep Notation formatting under \textbf{Notation.} blocks only.

\section{Error-based metrics}\label{appendix:metrics_error}
% MSE, RMSE, MAE, PSNR, PSNR-B, SNR, etc.
% Move full equations and variable definitions here.

\section{Similarity-based metrics}\label{appendix:metrics_similarity}
% SSIM, DSSIM, MS-SSIM, IW-SSIM, CW-SSIM, W-SSIM, ESSIM, G-SSIM, R-SSIM, L-SSIM, UQI, SCC, FSIM, FSIMc, SR-SIM, VSNR, etc.

\section{Information-theoretic metrics}\label{appendix:metrics_information}
% IFC, VIF, VIFp, VIFc.

\section{Deep feature-based metrics}\label{appendix:metrics_deep}
% LPIPS variants, DISTS, WaDIQaM.
% Include any backbone- or layer-selection notes.

\section{No-reference metrics}\label{appendix:metrics_nr}
% NIQE, PIQE, SER-FIQ, MagFace.
% Include preprocessing, feature assumptions, and any dataset-specific caveats.

\chapter{Regression models compared}\label{appendix:regressors}
This appendix summarizes alternative regressors evaluated for fusion and provides reference equations and training details.

\section{Linear and regularized linear models}\label{appendix:regressors_linear}
% OLS, Ridge, ElasticNet. Objective, penalties, cross-validation notes.

\section{Kernel methods}\label{appendix:regressors_kernel}
% SVR (RBF, linear). Margin, kernels, C and epsilon roles.

\section{Tree ensembles}\label{appendix:regressors_ensembles}
% Random Forest, Extra Trees. Bias-variance trade-offs, impurity measures.

\section{Gradient boosting families}\label{appendix:regressors_boosting}
% Gradient Boosting, HistGradientBoosting, XGBoost, LightGBM.
% Losses, regularization, histogram binning, sparsity handling.

% Note: CatBoost is documented in Methodology to emphasize it as the chosen model.

% =========================
% END SKELETON
% =========================
